% =========================================================
% Duality
% =========================================================

\section{Duality}
\label{sec:duality}

\begin{toolkitbox}{Toolkit: Duality}
Every partially ordered set has a dual obtained by reversing all order
relations.  Any statement that holds in every poset therefore holds in
every dual poset, which means its order-reversed version also holds
universally.  This is the Duality Principle — the formal engine behind
phrases like ``by duality, the corresponding result for infima follows.''
\end{toolkitbox}

% ---------------------------------------------------------
% def:dual-ordered-set
% ---------------------------------------------------------

\begin{definitionbox}{Definition (Dual Ordered Set)}
\begin{definition}[Dual Ordered Set]
\label{def:dual-ordered-set}
Let $(P, \leq)$ be a partially ordered set.  The \emph{dual} (or
\emph{opposite}) of $P$, denoted $P^{\mathrm{op}}$, is the set $P$
equipped with the reversed order:
\[
x \leq_{P^{\mathrm{op}}} y \;\Longleftrightarrow\; y \leq_P x.
\]
\end{definition}
\end{definitionbox}

\begin{remark*}[Standard quantified statement]
\[
x \leq_{P^{\mathrm{op}}} y
\;\coloneqq\;
y \leq_P x.
\]
\end{remark*}

\begin{remark*}[Interpretation]
Dualising a poset exchanges the roles of upper and lower: every upper bound
becomes a lower bound, every supremum becomes an infimum, and every maximal
element becomes a minimal element.  In a finite poset the Hasse diagram of
$P^{\mathrm{op}}$ is obtained from that of $P$ by turning it upside down.
The dual of a total order is again a total order.
\end{remark*}

\begin{remark*}[Dependencies]
\hyperref[def:lattice-order-partial-order]{Partial order}.
\end{remark*}

% ---------------------------------------------------------
% prop:duality-principle
% ---------------------------------------------------------

\begin{proposition}[Duality Principle]
\label{prop:duality-principle}
Let $\Phi$ be a statement about partially ordered sets, expressed purely
in terms of $\leq$, that holds in every partially ordered set.  Then the
dual statement $\Phi^{\mathrm{op}}$, obtained by replacing every
occurrence of $\leq$ with $\geq$, also holds in every partially ordered
set.
\end{proposition}

\begin{remark*}[Interpretation]
The Duality Principle eliminates the need to re-prove results whose dual
is already established.  Any theorem about upper bounds, suprema, maximal
elements, or top elements automatically yields a theorem about lower
bounds, infima, minimal elements, or bottom elements by applying the
principle to $P^{\mathrm{op}}$.
\end{remark*}

\begin{remark*}[Dependencies]
\hyperref[def:dual-ordered-set]{Dual ordered set}.
\end{remark*}

\begin{remark*}[Proof]
\hyperref[prf:duality-principle]{Go to proof.}
\end{remark*}
